\scanpage{511}
\mspage{558}

\[
\diagram{\begin{tikzcd}[ampersand replacement=\&,column sep=large]
L_0 \arrow[r] \& S\deleted{\unsure{\mathcal{T}}}\Gamma_{\rho,\sigma}^{0}\subset B_{\rho,\sigma}\times Z_\rho\times Z_\sigma \\
\ell \arrow[r,mapsto] \& (\ell,1,1)
\end{tikzcd}}\tag{40}
\]
and the action of $L_0$ by left translation on
$\unsure{S\mathcal{T}'_{\rho,\sigma}|S\Gamma'_{\rho,\sigma}}$
becomes, by $(39)$, the action on $S\unsure{\mathscr{G}|\mathcal{G}}_{\rho,\sigma}$ given by
\[
(\alpha,\beta)\longmapsto(\ell\alpha,\ell\beta).\tag{41}
\]

Let us therefore set
\margin{\unsure{denoted} $\unsure{\Gamma|\mathcal{T}}_{\rho,\sigma}^{0}$ in \S\,48, $\overline{\text{\uline{IX}}}$ $\unsure{\#}$.}
\[
\deleted{\unsure{\mathcal{T}}}\Gamma_{\rho,\sigma}^{0}
\overset{\mathrm{dfn}}{=}
S\deleted{\unsure{\mathcal{T}}}\Gamma_{\rho,\sigma}^{0}/\operatorname{Im}L_0
\simeq Z_\rho\mathbin{\times}_{\gamma_{\rho,\sigma}}Z_\sigma,
\tag{42}
\]
we find that $(39)$ induces a canonical bijection
\[
\diagram{\begin{tikzcd}[ampersand replacement=\&,column sep=large]
\deleted{\unsure{\mathcal{T}}}\Gamma_{\rho,\sigma}^{0}\arrow[r,"\sim"]\&L_0\backslash S\unsure{\mathscr{G}|\mathcal{G}}_{\rho,\sigma}
\end{tikzcd}}
\deleted{\unsure{\simeq Z_\rho\times Z_\sigma}}
\tag{43}
\]
\deleted{The expression $(37)$ for $S_0\mathcal{T}_{\rho,\sigma}$ gives an \inserted{exact} sequence,}
\[
\deleted{(44)\quad 1\longrightarrow L_0\times S}
\qquad\deleted{\text{Let us note the sequences}}
\]
Setting
\[
\left\{\begin{aligned}
\underset{\cdot}{S}Z_\rho&\overset{\mathrm{dfn}}{=}\operatorname{Ker}(Z_\rho\longrightarrow\Gamma_{\rho,\sigma})=\operatorname{Centr}_{\unsure{\mathcal{G}}}(\rho)\\
SZ_{\rho}&=\operatorname{Ker}(Z_{\rho}\mathbin{-}\Gamma_{\rho,\unsure{2|\sigma}})=\operatorname{Centr}_{\mathcal{G}}(\sigma)
\end{aligned}\right.\tag{44}
\]
\margin{also introduce
$S\Gamma_{\rho,\sigma}=\operatorname{Ker}(\Gamma_{\rho,\sigma}\longrightarrow\gamma_{\rho,\sigma}),$
$SG_{\rho,\sigma}=\operatorname{Ker}(G_{\rho,\sigma}\longrightarrow\gamma_{\rho,\sigma})=\mathcal{G}_{\rho,\sigma}=\mathcal{G}$.}

we thus have exact sequences
\[
\left\{\begin{gathered}
\diagram{\begin{tikzcd}[ampersand replacement=\&,column sep=small]
1\arrow[r,no head]\&SZ_\rho\arrow[r]\&Z_\rho\arrow[r,no head]\&\Gamma_{\unsure{\rho\sigma|\rho,\sigma}}\arrow[r,no head]\&1
\end{tikzcd}}\\
\diagram{\begin{tikzcd}[ampersand replacement=\&,column sep=small]
1\arrow[r,no head]\&SZ_\sigma\arrow[r,no head]\&Z_\sigma\arrow[r,no head]\&\Gamma_{\rho,\sigma}\arrow[r,no head]\&1
\end{tikzcd}}
\end{gathered}\right.\,,\tag{45}
\]
and the expressions $(37)$, $(42)$ for $S\deleted{\unsure{\mathcal{T}}}\Gamma_{\rho,\sigma}^{0}$, $\deleted{\unsure{\mathcal{T}}}\Gamma_{\rho,\sigma}^{0}$ give
\[
\diagram{\begin{tikzcd}[ampersand replacement=\&,column sep=small]
1\arrow[r]\&L_0\times SZ_\rho\times SZ_\sigma\arrow[r,no head]\&S_0\deleted{\unsure{\mathcal{T}}}\Gamma_{\rho,\sigma}^{0}\arrow[r,no head]\&\gamma_{\rho,\sigma}\arrow[r,no head]\&1
\end{tikzcd}}\tag{46}
\]
\[
\diagram{\begin{tikzcd}[ampersand replacement=\&,column sep=small]
1\arrow[r]\&SZ_\rho\times SZ_\sigma\arrow[r]\&\deleted{\unsure{S}}\deleted{\unsure{\mathcal{T}}}\Gamma_{\rho,\sigma}^{0}\arrow[r,no head]\&\gamma_{\rho,\sigma}\arrow[r]\&1
\end{tikzcd}}.\tag{47}
\]
\uline{N.B.} In the case $\mathcal{G}=\mathrm{Sl}(2,\widehat{\mathbb{Z}})$, $\rho$ and $\sigma$ ``modular'', we recover \deleted{\illegible{1}} the $S_0\deleted{\unsure{\mathcal{T}}}\Gamma_{\rho,\sigma}^{0\unsure{4|\#}}$ and $\deleted{\unsure{\mathcal{T}}}\Gamma_{\rho,\sigma}^{0}$ of \S\,48 $\overline{\text{IX}}$.

\mspage{510}

\scanpage{512}
\mspage{559}

Now let
\[
\begin{gathered}
S_0M_{\rho,\sigma}^{0*}\subset
S\deleted{\unsure{\mathcal{T}}}\Gamma_{\rho,\sigma}^{0}\times G_{\rho,\sigma}
\quad(\subset B_{\rho,\sigma}\times Z_\rho\times Z_\sigma\times G_{\rho,\sigma})\\
\overset{\mathrm{dfn}}{=}\\
\left\{(u,a,b,U)\,\middle|\,
\underbrace{\delta_B(u)=\delta_\rho(a)=\delta_\sigma(b)=\delta_G(U)}_{\text{relation in }\gamma_{\rho,\sigma}\ \text{cf. }(17)}\right\}
\end{gathered}\tag{49}
\]
\[
\begin{gathered}
M_{\rho,\sigma}^{0}\subset T_\rho\times T_\sigma\times G_{\rho,\sigma}\\
\overset{\mathrm{dfn}}{=}\\
\left\{(a,b,U)\,\middle|\,
\underbrace{\delta_\rho(a)=\delta_\sigma(b)=\delta_G(U)}_{\text{relations in }\gamma_{\rho,\sigma}}\right\}
\end{gathered}\tag{50}
\]
We thus have a canonical homomorphism
\[
\diagram{\begin{tikzcd}[ampersand replacement=\&,column sep=large]
SM_{\rho,\sigma}^{0}\arrow[r]\&M_{\rho,\sigma}^{0}\\
(u,a,b,U)\arrow[r]\&\deleted{\unsure{(u}}(a,b,U)
\end{tikzcd}}\tag{51}
\]
whose kernel consists of the $(u,1,1,1)$ where $u\in B_{\rho,\sigma}$ is such that $\delta_B(u)=1$ i.e. $u\in L_0$, whence an exact sequence
\[
\diagram{\begin{tikzcd}[ampersand replacement=\&,column sep=small]
1\arrow[r,no head]\&L_0\arrow[r]\&SM_{\rho,\sigma}^{0}\arrow[r]\&M_{\rho,\sigma}^{0}\arrow[r,no head]\&1
\end{tikzcd}}\tag{52}
\]
and canonical homomorphisms
\[
\diagram{\begin{tikzcd}[ampersand replacement=\&,column sep=large,row sep=small]
\&G_{\rho,\sigma}\\
M_{\rho,\sigma}^{0}\arrow[ur,"\Theta_G"]\arrow[r,"\Theta_\rho",very thick]\arrow[dr,"\Theta_\sigma"']\&Z_\rho\\
\&Z_\sigma
\end{tikzcd}}\tag{53}
\]
\[
\left(\text{resp.}\quad
\diagram{\begin{tikzcd}[ampersand replacement=\&,column sep=small,row sep=small]
\&U\\
(a,b,U)\arrow[ur,mapsto]\arrow[r,mapsto]\arrow[dr,mapsto]\&a\\
\&b
\end{tikzcd}}\right),
\]
\unsure{giving back} the analogous homomorphisms on $SM_{\rho,\sigma}^{0}$ by composition with $(41)$, and moreover we have a canonical homomorphism
\[
\diagram{\begin{tikzcd}[ampersand replacement=\&,column sep=large]
SM_{\rho,\sigma}^{0}\arrow[r,"\Theta_B"]\&B_{\rho,\sigma}\\
(u,a,b,U)\arrow[r,mapsto]\&u
\end{tikzcd}}.\tag{54}
\]

\mspage{511}

\scanpage{513}

We also have \deleted{\unsure{of}} the interpretation
\[
S_0M_{\rho,\sigma}^{0}\simeq M_{\rho,\sigma}^{0}\mathbin{\times}_{\gamma_{\rho,\sigma}}B_{\rho,\sigma};\tag{55}
\]
( $S_0M_{\rho,\sigma}^{0}$ is expressed in terms of $M_{\rho,\sigma}^{0}$ and $B_{\rho,\sigma}$, and we \unsure{see} that the extension of $M_{\rho,\sigma}$ by $L_0$ is the inverse image of the extension $B_{\rho,\sigma}$ of $\gamma_{\rho,\sigma}$ by $L_0$. And we have
\[
M_{\rho,\sigma}^{0}\simeq
\underbrace{Z_\rho\mathbin{\times}_{\gamma_{\rho,\sigma}}Z_\sigma}_{\deleted{\unsure{\mathcal{T}}}\Gamma_{\rho,\sigma}^{0}}
\mathbin{\times}_{\gamma_{\rho,\sigma}}G_{\rho,\sigma}.\tag{56}
\]
whence we deduce (since $\delta_\rho,\delta_\sigma,\delta_G$ are epimorphisms)
\[
\diagram{\begin{tikzcd}[ampersand replacement=\&,column sep=small]
1\arrow[r,very thick]\&\deleted{\unsure{\Gamma_\rho|T_\rho}}\underbrace{SZ_\rho\times SZ_\sigma\times\unsure{\mathcal{G}}}_{\overset{\mathrm{dfn}}{=}M_{\rho,\sigma}^{0+}}\arrow[r]\&\unsure{M_{\rho,\sigma}'|M_{\rho,\sigma}^{0}}\arrow[r,"\delta_M^{\unsure{\#|*}}"]\&\gamma_{\rho,\sigma}\arrow[r]\&1
\end{tikzcd}}\tag{57}
\]
where we have \deleted{\unsure{set}} $(44)$
\[
\deleted{(\unsure{48|58})}\qquad
\left\{\begin{aligned}
SZ_\rho&=Z_\rho\cap\mathcal{G}={\underset{\deleted{\unsure{\mathrm{aut}}}}{\uline{\operatorname{Centr}}}}_{\mathcal{G}}(\rho)\\
SZ_\sigma&=Z_\sigma\cap\mathcal{G}={\underset{\deleted{\unsure{\mathrm{aut}}}}{\uline{\operatorname{Centr}}}}_{\mathcal{G}}(\sigma)
\end{aligned}\right.
\]
Similarly we have
\[
S_0M_{\rho,\sigma}^{0}\simeq
\underbrace{B_{\rho,\sigma}\mathbin{\times}_{\gamma_{\rho,\sigma}}Z_\rho\mathbin{\times}_{\gamma_{\rho,\sigma}}Z_\sigma}_{S_0\deleted{\unsure{\mathcal{T}}}\Gamma_{\rho,\sigma}^{0}}
\mathbin{\times}_{\gamma_{\rho,\sigma}}G_{\rho,\sigma}
\tag{49}
\]
whence we have a canonical exact sequence
\[
\diagram{\begin{tikzcd}[ampersand replacement=\&,column sep=small]
1\arrow[r]\&L_0\times SZ_\rho\times SZ_\sigma\times\mathcal{G}\arrow[r]\&S_0M_{\rho,\sigma}^{0}\arrow[r]\&\Gamma_{\rho,\sigma}\arrow[r,no head]\&1
\end{tikzcd}}\tag{50}
\]
Let us consider the canonical homomorphism
\[
\diagram{\begin{tikzcd}[ampersand replacement=\&,column sep=large]
\mathcal{G}\arrow[r,"i_{\mathcal{G}}"]\&M_{\rho,\sigma}^{0}\\
U\arrow[r,mapsto]\&(1,1,U)
\end{tikzcd}},\tag{51}
\]

\mspage{512}

\scanpage{514}
\mspage{560}

its image is the third factor in the kernel of $(57)$, and we have a canonical exact sequence (deduced from $(56)$--$(57)$)
\[
\diagram{\begin{tikzcd}[ampersand replacement=\&,column sep=small]
1\arrow[r]\&\underset{\inserted{\substack{\text{sometimes denoted}\\\mathcal{G}_{\rho,\sigma}?}}}{\unsure{\mathcal{G}}}\arrow[r]\&M_{\rho,\sigma}^{0}\arrow[r]\&\deleted{\unsure{\mathcal{T}}}\Gamma_{\rho,\sigma}^{0}\arrow[r,no head]\&1
\end{tikzcd}}\tag{52}
\]
\margin{$(53')$ \diagram{\begin{tikzcd}[ampersand replacement=\&,column sep=small]
1\arrow[r]\&SZ_\rho\times SZ_\sigma\arrow[r]\&M_{\rho,\sigma}^{0}\arrow[r]\&\unsure{G|G_1}\arrow[r,no head]\&1
\end{tikzcd}}}
and similarly we have (via $(49)$ or $50)$)
\[
\left\{\begin{aligned}
&(53)\quad\diagram{\begin{tikzcd}[ampersand replacement=\&,column sep=small]
1\arrow[r]\&\mathcal{G}\arrow[r]\&S_0M_{\rho,\sigma}^{0}\arrow[r]\&S_0\deleted{\unsure{\mathcal{T}}}\Gamma_{\rho,\sigma}^{0}\arrow[r]\&1
\end{tikzcd}}.\\[1ex]
&\text{Let us consider the map}\\[1ex]
&(54)\quad\diagram{\begin{tikzcd}[ampersand replacement=\&,column sep=large]
S_0M_{\rho,\sigma}^{0}\arrow[r]\&\mathcal{G}\times\mathcal{G}\times\mathcal{G}\\
\{u,a,b,U\}\arrow[r,mapsto]\&(ua^{-1},ub^{-1},uU^{-1})=(\alpha,\beta,f)
\end{tikzcd}}
\end{aligned}\right.
\]
we \unsure{see} \deleted{\illegible{1}} by the corollary on page\,557 (cf. $(37)$) that this induces a bijection
\[
S_0M_{\rho,\sigma}^{0}\simeq(S\unsure{\mathscr{G}|\mathcal{G}}_{\rho,\sigma})\times\mathcal{G}.
\tag{55}
\]
We thus have, for an $x\in S_0M_{\rho,\sigma}^{0}$, on the one hand the quantities
\[
\diagram{\begin{array}{cccc}
u,&a,&b,&U\\
\rotatebox{90}{$\in$}&\rotatebox{90}{$\in$}&\rotatebox{90}{$\in$}&\rotatebox{90}{$\in$}\\
B_{\rho,\sigma}&Z_\rho&Z_\sigma&G_{\rho,\sigma}
\end{array}},\tag{56}
\]
associated, on the other hand the elements
\[
\alpha,\beta,f,\quad\unsure{\zeta|\xi},\eta,\quad\alpha_0\ \beta_0,g,h
\qquad\text{all }\in\mathcal{G}\tag{57}
\]
\inserted{where $(\alpha,\beta,f)$} defined in terms of $(56)$ in $(54)$, and where we have set
\[
\left\{\begin{aligned}
\zeta&=\alpha\rho(\alpha)^{-1},\quad
\eta=\beta\sigma(\beta)^{-1},\quad
\alpha_0=f^{-1}\alpha\unsure{\#|\zeta},\quad\beta_0=f^{-1}\beta,\\
g&=\alpha_0\rho(\alpha_0)^{-1}=f^{-1}\zeta\rho(f),\quad
h=\beta_0\sigma(\beta_0)^{-1}=f^{-1}\eta\sigma(f)
\end{aligned}\right.\tag{58}
\]

\mspage{513}

\scanpage{515}

We will have, in terms of $\alpha,\beta,f$, the following expressions for the action of $U$ on $\mathcal{G}$:
\[
\left\{\begin{aligned}
U(\rho)&=\operatorname{int}(f^{-1}\alpha)\deleted{\unsure{\operatorname{int}(\alpha)}}\rho
 =\operatorname{int}(f^{-1})(\zeta\rho)\\
U(\sigma)&=\operatorname{int}(f^{-1}\beta)\deleted{\unsure{\operatorname{int}(\beta)}}(\sigma)
 =\operatorname{int}(f^{-1})(\eta\sigma).
\end{aligned}\right.\tag{59}
\]
\deleted{\unsure{and}\ $\unsure{\zeta|\xi},\unsure{\eta}$\unsure{.}}

\uline{\unsure{Remark|Remark}} Finally, in all these developments, we have not really used the so-called ``fundamental hypothesis'' $(20)\deleted{\unsure{b}}$, \deleted{\unsure{and}} except part a). In the absence of b), one should introduce
\[
\begin{gathered}
\Sigma_{\rho,\sigma}^{\mathrm{eff}}\subset\Sigma_{\rho,\sigma}\subset\mathcal{G}\times\mathcal{G}\\[-1pt]
\begin{array}{c@{\qquad}l}
\rotatebox{90}{$=$}&\\[-3pt]
\multicolumn{2}{l}{\left\{(\zeta,\eta)\in\unsure{\mathcal{G}|\Sigma_\rho}\times\mathcal{G}\mid
\exists u\in B_{\rho,\sigma},\ u(\rho)=\zeta\rho,\ u(\sigma)=\eta\sigma\right\},}
\end{array}
\end{gathered}\tag{60}
\]
which is thus nothing other than the image of the map $(9)$: \uline{effective} pairs $(\zeta,\eta)$ satisfying the conditions defining \deleted{\unsure{the}\ $\unsure{(5)}$}\inserted{$\Sigma_{\rho,\sigma}$ (cf. $(8)$)}\unsure{)}

\deleted{The map
$\diagram{\begin{tikzcd}[ampersand replacement=\&,column sep=large]
S\unsure{\mathscr{G}|\mathcal{G}}_{\rho,\sigma}\arrow[r]\&\Sigma_{\rho,\sigma}
\end{tikzcd}}$
\inserted{$(28)$} (will not necessarily be surjective, its image \unsure{is|will be} $\Sigma_{\rho,\sigma}^{\mathrm{eff}}\inserted{\simeq B_{\rho,\sigma}}$), otherwise there is nothing to change.}
\margin{\unsure{then|by}\par\unsure{bijection|bis}}

let us write
\[
\begin{aligned}
S\unsure{\mathscr{G}|\mathcal{G}}_{\rho,\sigma}^{\mathrm{eff}}
 &=\text{Inverse image of }\Sigma_{\rho,\sigma}^{\mathrm{eff}}\\
 &\text{under }\diagram{\begin{tikzcd}[ampersand replacement=\&,column sep=large]
S\unsure{\mathscr{G}|\mathcal{G}}_{\rho,\sigma}\arrow[r,no head]\&\Sigma_{\rho,\sigma}
\end{tikzcd}}\quad(28),
\end{aligned}\tag{61}
\]
\deleted{the inverse image}
and it is on $S\unsure{\mathscr{G}|\mathcal{G}}_{\rho,\sigma}^{\mathrm{eff}}$ that we will have a structure

\mspage{514}

\scanpage{516}
\mspage{561}

of a group \unsure{isomorphic} \unsure{to|:} $S_0\mathcal{T}_{\rho,\sigma}^{0}$, by the isomorphism (which replaces (38))
\diagram{\begin{tikzcd}[ampersand replacement=\&,column sep=large]
 S_0\deleted{\unsure{\mathcal{T}}}\Gamma_{\rho,\sigma}^{0} \arrow[r,phantom,"\simeq"] \& S_0\unsure{\mathscr{G}|\mathcal{G}}_{\rho,\sigma}^{\mathrm{eff}}\,. \\
 (\underline{u},\underline{a},\underline{b}) \arrow[r] \& (\underbrace{ua^{-1}}_{\zeta},\underbrace{ub^{-1}}_{\eta})
\end{tikzcd}\quad(62)}

This being so, let us return to $M_{\rho,\sigma}^{\sim}$, and the exact sequence
\diagram{\begin{tikzcd}[ampersand replacement=\&,column sep=small]
 1 \arrow[r] \& \mathcal{G}_{0,3}^{+\wedge} \arrow[r] \& M_{\unsure{0,3|\rho,\sigma}}^{\sim} \arrow[r] \& \underbrace{\Gamma_{0,3}^{\sim}}_{\Gamma_{\mathbb Q}^{\sim}} \arrow[r] \& 1
\end{tikzcd}\quad(63).}

Consider the subgroup
\[
\begin{aligned}
(64)\qquad
\Gamma_{\mathbb Q}^{\prime\sim}
 &= \operatorname{Ker}\bigl(\Gamma_{0,3}^{\sim}
 \xrightarrow{(\varepsilon_3,\deleted{\unsure{\varepsilon_2}})}
 \{\pm1\}\deleted{\unsure{\times\{\pm1\}}}\bigr)\\
 &= \{\gamma\in\Gamma_{0,3}^{\sim}\mid\mu(\gamma)\equiv1\ (\unsure{3}\deleted{\unsure{12}})\}.
\end{aligned}
\]
\deleted{\illegible{3}}
\deleted{\unsure{two factors}}

and let \deleted{$\unsure{M_{0,3}^{\prime\sim}}$}
\[
 (65)\qquad M_{0,3}^{\prime\sim}
 = \text{\unsure{Inverse image} of }\Gamma_{0,3}^{\prime\sim}\text{ in }M_{0,3}^{\sim}.
\]

Let $\mathcal{G}$ be a profinite group, and
\[
 (66)\qquad\psi:\mathcal{G}_{0,3}^{+\wedge}\longrightarrow\mathcal{G},
\]
a \inserted{continuous} surjective homomorphism --- which amounts to the same thing as giving two elements
\[
 (67)\qquad\rho_{\mathcal{G}},\unsure{\sigma_{\mathcal{G}}}\in\mathcal{G}
\]
(the images of $\rho,\sigma$ under $\psi$) satisfying
\[
 (68)\qquad\rho_{\mathcal{G}}^3=\sigma_{\mathcal{G}}^2=1.
\]
\mspage{515}

\scanpage{517}

\margin{introduce it by equivalent conditions:

\unsure{a)}
\diagram{\begin{tikzcd}[ampersand replacement=\&,column sep=large]
 M_{0,3}^{\prime\sim}[0] \arrow[r] \& \Sigma_{\rho,\sigma}^{\text{\textcircled{eff}}}
\end{tikzcd}} \unsure{all}

b) $\forall u\in M_{0,3}^{\prime\sim},\ \exists u_{\unsure{\mathcal{G}}}\in\operatorname{Aut}\mathcal{G}$ such that
$u_{\unsure{M|\mathcal{G}}}\circ\psi=\psi\circ u$

b') Idem with $u\in M_{0,3}^{\prime\sim}[0]$

\unsure{c)} $\exists\psi_M$ which extends $\psi$

c') $\exists\psi_G$ which extends $\psi$

We will therefore say that \deleted{\unsure{effective}} \unsure{automatic|automorphism}

$\psi$ is \uline{\unsure{effective}} if (20) b) satisfied}

\uline{Definition} We will say that $\psi$ is \uline{effective} if for every
$u\in M_{0,3}^{\prime\sim}[0]$, putting $u(\rho)=\zeta\rho$, $u(\sigma)=\eta\sigma$,
\[
 (\zeta_{\mathcal{G}},\eta_{\mathcal{G}})
 =(u(\zeta),u(\eta))\in\Sigma_{\rho,\sigma}
\]
effective (i.e. $\in\Sigma_{\rho,\sigma}^{\mathrm{eff}}$, i.e. $\exists u_{\mathcal{G}}$ \unsure{automorphism} $\mathcal{G}$ such that we have \deleted{$\unsure{u_{\mathcal{G}}\rho}$}
\[
 (69)\qquad
 u_{\mathcal{G}}(\rho_{\mathcal{G}})=\zeta_{\mathcal{G}}\rho_{\mathcal{G}},
 \qquad
 u_{\mathcal{G}}(\sigma_{\mathcal{G}})=\eta_{\mathcal{G}}\sigma_{\mathcal{G}}
\]
or, what amounts to the same thing, that we have commutativity in
\diagram{\begin{tikzcd}[ampersand replacement=\&,column sep=large,row sep=large]
 \mathcal{G}_{0,3}^{+\wedge} \arrow[r,"u"] \arrow[d,"\psi"'] \& \mathcal{G}_{0,3}^{+\wedge} \arrow[d,"\psi"] \\
 \mathcal{G} \arrow[r,"u_{\mathcal{G}}"'] \& \mathcal{G}
\end{tikzcd}\quad(70).}

(Thus to say that this is the case for \unsure{every} $u\in M_{0,3}^{\prime\sim}\deleted{[0]}$ means that these $u$ transform $\operatorname{Ker}\psi$ into itself --- or also, that
$M_{0,3}^{\prime\sim}\deleted{\unsure{[0]}}$
$(=M_{0,3}^{\prime\sim}[0].\mathcal{G}_{0,3}^{+\wedge})$ in its entirety stabilizes $\operatorname{Ker}\psi$.
\deleted{When} This condition is therefore automatically fulfilled when $\mathcal{G},\rho_{\mathcal{G}},\sigma_{\mathcal{G}}$ satisfy \unsure{the condition|the conditions} b) of (20).

We have done what was needed to have \unsure{the}
\mspage{516}

\scanpage{518}
\mspage{562}

\uline{Theorem} Let $\psi=\psi_{\mathcal{G}}:\mathcal{G}_{0,3}^{+\wedge}\longrightarrow\mathcal{G}$ be an effective surjective homomorphism (cf. above).
Then there exists a unique homomorphism
\[
 (71)\qquad\psi_M:\deleted{\unsure{S_0}}M_{0,3}^{\prime\sim\deleted{\unsure{*}}}
 \longrightarrow\deleted{\unsure{S_0}}M_{\rho_{\mathcal{G}},\sigma_{\mathcal{G}}}^{0}
\]
\deleted{making the diagram commutative}
\diagram{\deleted{\begin{tikzcd}[ampersand replacement=\&,column sep=large,row sep=large]
 \mathcal{G}_{0,3}^{+\wedge} \arrow[r,"\unsure{\psi|\varphi}"] \arrow[d,"\psi"'] \& \deleted{\unsure{\mathcal{G}}}\,M_{0,3}^{\prime\sim} \arrow[d,"\psi"] \\
 \deleted{\unsure{\mathcal{G}}}\,\mathcal{G} \arrow[r] \& M_{\rho_{\mathcal{G}},\sigma_{\mathcal{G}}}^{0}\ (\overset{\mathrm{dfn}}{=}M_{\psi}^{0})
\end{tikzcd}\quad(72)}}

\margin{In fact, we have \unsure{of course in} the first place a homomorphism
\diagram{\begin{tikzcd}[ampersand replacement=\&,column sep=large]
 S_0M_{0,3}^{\prime\sim} \arrow[r] \& S_0M_{\psi}^{0}.
\end{tikzcd}}}

having the following property: if \deleted{$\unsure{u\in}$}
$u\unsure{=|\in}M_{0,3}^{\prime\sim}$ is of the form $u_{\alpha,\beta,f}$
($\alpha,\beta,f\in\mathcal{G}_{0,3}^{+\wedge}$, satisfying the loop equation
$\alpha\rho(\alpha)^{-1}=\beta\sigma(\beta)^{-1}\ell_1^{\nu}$), then
$\psi_M(u)$ is the element of \deleted{$\unsure{M_{\rho_{\mathcal{G}},\sigma_{\mathcal{G}}}^{0}}$} $M_{\psi}^{0}$ defined by the invariants $\psi(\alpha),\psi(\beta);\psi(f)$.

\uline{Corollary} We then have homomorphisms of exact sequences
\diagram{\begin{tikzcd}[ampersand replacement=\&,column sep=small,row sep=large]
 1 \arrow[r] \& \mathcal{G}_{0,3}^{+\wedge} \arrow[r,no head] \arrow[d,"\psi_{\mathcal{G}}"'] \& M_{0,3}^{\prime\sim} \arrow[r] \arrow[d,"\psi_M"] \& \Gamma_{\mathbb Q}^{\prime\sim} \arrow[r] \arrow[d,"\unsure{\psi_{\mathcal{T}}|\psi_{\Gamma}}"] \& 1 \\
 1 \arrow[r] \& \mathcal{G} \arrow[r] \& M_{\psi}^{0} \arrow[r] \& \deleted{\unsure{\mathcal{T}}}\Gamma_{\psi}^{0} \arrow[r] \& 1
\end{tikzcd}\quad(72).}

(The homomorphism $\psi_M$ is made explicit by means of three homomorphisms
\mspage{517}

\scanpage{519}

\[
\left\{\begin{aligned}
\psi_{T_\rho}&:\Gamma_{\mathbb Q}^{\prime\sim}\longrightarrow \unsure{Z_\rho|T_\rho}\\
\psi_{T_\sigma}&:\Gamma_{\mathbb Q}^{\sim}\longrightarrow \unsure{Z_\sigma|T_\sigma}
\end{aligned}\right.\qquad(73)
\]
\[
\psi_G:M_{0,3}^{\prime\sim}\longrightarrow G_{\rho,\sigma}\qquad(74)
\]
\margin{N.B. for simplicity we write $\rho,\sigma$ instead of $\rho_{\mathcal G},\sigma_{\mathcal G}$ as subscripts}

\deleted{satisfying the conditions}
(For \deleted{$\unsure{u}$} $x\in M_{0,3}^{\prime\sim}$,
\inserted{$\dot x$ its image in $\Gamma_{\mathbb Q}^{\prime\sim}$,}
\deleted{$a=\psi_{\unsure{T_\rho|\Gamma_\rho}}(\dot x),\ b=\psi_{\unsure{T_\rho|\Gamma_\rho}}(\unsure{y|\dot x}),\ U=\unsure{\psi u|\psi(x)}$}
\[
\deleted{\unsure{\delta_G(}}\quad
 a=\psi_{\unsure{T_\rho|\Gamma_\rho}}(\dot x),\quad b=\psi_{\unsure{T_\sigma|\Gamma_\sigma}}(\dot x),\quad U=\psi_G(x),
\qquad(75)
\]
we have
\[
\delta_\rho(a)=\delta_\sigma(b)=\delta_G(U)\quad\text{in }\gamma_{\rho,\sigma}\qquad(76)
\]
i.e. the three composite homomorphisms
\diagram{\begin{tikzcd}[ampersand replacement=\&,column sep=large,row sep=large]
\&\Gamma_{\mathbb Q}^{\prime\sim}\arrow[r,"\psi_{T_\rho}"]\&T_\rho\arrow[rd,"\delta_\rho"]\&\\
M_{0,3}^{\prime\sim}\arrow[ru,"\mathrm{can}"]\arrow[r,"\mathrm{can}"]\arrow[rrd,"\psi_G"',"\unsure{1}"{pos=.43}]\arrow[rrd,shorten >=55pt,shorten <=40pt]\&\Gamma_{\mathbb Q}^{\prime\sim}\arrow[r,no head,"\psi_{T_\sigma}"]\&T_\sigma\arrow[r,"\delta_\sigma"]\&\gamma_{\rho,\sigma}\arrow[ld,"\delta_G"]\\
\&\&G_{\rho,\sigma}\&
\end{tikzcd}\qquad(77)}
are equal. Moreover, we will have by (71),
\[
\psi_G|\mathcal G_{0,3}^{+\wedge}=\psi_{\mathcal G},\qquad(78)
\]
and finally this:

(79) If $x\in \deleted{\unsure{M}}M_{0,3}^{\prime\sim}$, $U=\psi_G(x)$, we have commutativity in
\diagram{\begin{tikzcd}[ampersand replacement=\&,column sep=huge,row sep=large]
\mathcal G_{0,3}^{+\wedge}\arrow[r,"\psi"]\arrow[d,"\operatorname{int}(x)|\mathcal G_{0,3}^{+\wedge}"']\&\mathcal G\arrow[d,"\operatorname{int}(U)|\mathcal G"]\\
\mathcal G_{0,3}^{+\wedge}\arrow[r,"\psi"]\&\mathcal G
\end{tikzcd}\qquad(80)}

\hfill $.\ .\ .\ .\ .\ .$ 
\mspage{518}

\scanpage{520}
\mspage{563}

. We can regard all this as a statement of \uline{functoriality} of the constructions made, with respect to a surjective homomorphism
\diagram{\begin{tikzcd}[ampersand replacement=\&]
\mathcal G\arrow[r,"\varphi=\varphi_{\mathcal G}"]\&\mathcal G'
\end{tikzcd}\qquad(81)}
\margin{\unsure{It would have} been better to take \unsure{then}
\diagram{\begin{tikzcd}[ampersand replacement=\&]\mathcal G'\arrow[r,"\varphi"]\&\mathcal G\end{tikzcd}},
in order to be able to specialize to $\mathcal G'$}
transforming the pair $(\rho,\sigma)$ of topological generators of $\mathcal G$ into a pair $(\rho',\sigma')$.
There is then a unique homomorphism
\diagram{\begin{tikzcd}[ampersand replacement=\&]
S_0M_{\mathcal G}^{0}\arrow[r,"\varphi_{M^0}"]\&S_0M_{\mathcal G'}^{0}
\end{tikzcd}\qquad(82)}
satisfying the following conditions:

\margin{this shows that $\varphi_{SM}$ is determined when we know the four homomorphisms $\varphi_B,\varphi_\rho,\varphi_\sigma,\varphi_G$}
a) $\unsure{\varphi_M|\varphi_{SM}}$ \unsure{passes} to quotients \unsure{as|to} homomorphisms
\[
\left\{\begin{aligned}
\varphi_B&:B_{\rho,\sigma}\longrightarrow\unsure{Z_{\rho,\sigma}|\Sigma_{\rho,\sigma}}\\
\varphi_\rho&:Z_\rho\longrightarrow Z_{\rho'}\\
\varphi_\sigma&:Z_\sigma\mathrel{\relbar\joinrel\relbar} Z_{\sigma'}\\
\varphi_G&:G_{\rho,\sigma}\mathrel{\relbar\joinrel\relbar} G_{\rho',\sigma'}
\end{aligned}\right.\ ;\qquad(83)
\]

b) \deleted{We have commutativity in the diagram,}
\inserted{\deleted{The homomorphisms $\varphi_B,\varphi_\rho,\varphi_\sigma$ \unsure{are induced by}}}
\[
(84)\quad\left\{\begin{gathered}
\deleted{\text{\unsure{only}}}\quad\varphi_G\qquad
\text{(which shows that knowledge}\\
\text{\illegible{1} of }\varphi_{SM}\ \unsure{b|se}\ \text{amounts to that of }\varphi_G\text{)};
\end{gathered}\right.
\]

c) We have commutativity in
\diagram{\begin{tikzcd}[ampersand replacement=\&,column sep=large,row sep=large]
\mathcal G\arrow[r]\arrow[d,"\varphi=\varphi_{\mathcal G}"']\&G_{\rho,\sigma}\arrow[d,"\varphi_G"]\\
\mathcal G'\arrow[r]\&G_{\rho',\sigma'}
\end{tikzcd}\qquad(85)}
(which shows, since $G_{\rho,\sigma}=B_{\rho,\sigma}\cdot\mathcal G$, that $\varphi_G$ is known when we know $\varphi_B$;

d) For $u\in B_{\rho,\sigma}$, \deleted{\unsure{denoting} by $u'$} \inserted{putting $u'=\varphi_B(u)$}, we have commutativity in
\diagram{\begin{tikzcd}[ampersand replacement=\&,column sep=large,row sep=large]
\mathcal G\arrow[r,"u"]\arrow[d,"\varphi_{\mathcal G}"']\&\mathcal G\arrow[d,"\varphi_{\mathcal G}"]\\
\mathcal G'\arrow[r,"u'"]\&\mathcal G'
\end{tikzcd}\qquad(86)}
\mspage{519}

\scanpage{521}

-- (which determines $\varphi_B$ in terms of $\varphi=\varphi_{\mathcal G}$, thus \unsure{this|it} determines $\unsure{\varphi_{S_0M}}$, and by passage to the quotient, also
\[
 (87)\qquad \varphi_M:M_{\rho,\sigma}\longrightarrow M_{\rho',\sigma'}.
\]

Another way, more direct, of describing $\varphi_{S_0M^0}$ or $\varphi_M$, in terms of the \deleted{coordinates} \inserted{``parameters''} $(\alpha,\beta,f)$ instead of the $(u,a,b,U)$, is to say that one \deleted{$\unsure{\varphi_{S_0M^0}}$ \unsure{transforms}} transforms a $u_{\alpha,\beta,f}$ into the $u_{\alpha',\beta',f'}$, where $\alpha',\beta',f'$ are the images of $\alpha,\beta,f$ under $\varphi$:
\[
 (88)\qquad
 \varphi_{\deleted{\unsure{\mathcal G}}S_0\unsure{M_0|M^0}}
 \left(\underbrace{u_{\alpha,\beta,f}}_{\in\,S_0M^0_{\rho,\sigma}}\right)
 =u_{\alpha',\beta',f'},\quad
 \left\{\begin{aligned}
 \alpha'&=\varphi(\alpha),\\
 \beta'&=\varphi(\beta),\\
 f'&=\varphi(f).
 \end{aligned}\right.
\]

We thus recover, in a slightly more precise form, the preceding \unsure{theorem}, by \deleted{\unsure{putting}} applying it to the case where $\mathcal G=\mathcal G_{0,3}^{+\wedge}$ (the ``universal case'')$^*$, $\mathcal G'$ arbitrary --- it \unsure{is} better in the notation (81) to interchange the role of $\mathcal G$ and $\mathcal G'$.
\margin{$^*$ at least ``universal'' with respect to the \unsure{additional} hypothesis $\rho^3=1,\ \sigma^2=1$ in $\mathcal G$.}
It is good to \deleted{\unsure{make explicit}}\inserted{make explicit the essential invariants} (in the ``universal'' case:

\deleted{for}
\[
\begin{aligned}
 (90)\qquad
 \mathcal G'&=\mathcal G_{0,3}^{+\wedge},\qquad
 B_{\rho',\sigma'}^{\deleted{\unsure{0}}\,\prime}
 =M_{0,3}^{\prime\sim}[0],\qquad
 \gamma_{\rho',\sigma'}^{\unsure{\prime|0}}
 =\Gamma_{\mathbb Q}^{\prime\sim},\\
 G_{\rho',\sigma'}^{\unsure{\prime|0}}
 &=M_{0,3}^{\prime\deleted{\unsure{0}}\sim},\qquad
 Z_{\rho}^{\prime\deleted{\unsure{0}}}
 =M_{0,3}^{\prime\sim}\unsure{\left(-\frac13\right)|\left(\frac13\right)},\qquad
 Z_{\sigma}^{\prime\deleted{\unsure{0}}}
 =M_{0,3}^{\prime\sim}\unsure{\left(\frac12\right)|\left(-\frac12\right)},\\
 \Gamma_{\rho',\sigma'}^{\prime}&=\Gamma_{\mathbb Q}^{\prime\sim}
 \quad\deleted{\left(\text{hence}\quad
 Z_{\rho}^{\prime}\xrightarrow{\sim}\gamma_{\rho',\sigma'}^{\unsure{\prime|0}},\quad
 Z_{\sigma}^{\prime}\xrightarrow{\sim}\gamma_{\rho',\sigma'}^{\prime}\right)}.
\end{aligned}
\]
\inserted{hence}
\mspage{520}


\scanpage{522}
\mspage{564}
\[
\begin{gathered}
\Gamma_{\rho',\sigma'}^{\prime}\xrightarrow{\sim}\gamma_{\rho',\sigma'}^{\prime},\quad
Z_{\rho'}^{\prime}\xrightarrow{\sim}\gamma_{\rho',\sigma'}^{\prime},\quad
Z_{\sigma'}^{\prime}\xrightarrow{\sim}\gamma_{\rho',\sigma'}^{\prime},\quad\text{thus}\\
\unsure{(91)|(92)}\qquad
S_0M_{\rho',\sigma'}^{\prime}\simeq
\overset{\inserted{S_0G_{\rho',\sigma'}^{\prime}\unsure{\simeq|=}}}{S_0M_{0,3}^{\prime\sim}},\quad
\deleted{\unsure{G_{\rho',\sigma'}=}}M_{\rho',\sigma'}^{\prime}\xrightarrow{\sim}
\overset{\inserted{G_{\rho',\sigma'}^{\prime}\xrightarrow{\sim}}}{\deleted{\unsure{G}}M_{0,3}^{\prime\sim}}.
\end{gathered}
\]
\margin{hence
\[
(92)\quad\left\{\begin{aligned}
SZ_{\rho'}^{\prime}&=SZ_{\sigma'}^{\prime}\\
&=S\unsure{\Gamma|\mathcal T}_{\rho',\sigma'}^{\prime}=1
\end{aligned}\right.
\quad(\unsure{cf.})\ \text{of}\ (92)
\]}

\uline{Remark} In writing the relations (86), it became evident that the notations used in this section II (from page\ \unsure{0}\ 553') were inadequate, and that they should be modified as follows:
\margin{\unsure{This had been suggested to me in the case} $\mathcal G=\mathrm{Sl}(2,\widehat{\mathbb Z})$, \unsure{modular pair} $\rho,\sigma$ \unsure{which does not come from} $\mathcal G_{0,3}^{+\wedge}$, because of the $\{\pm1\}$! \unsure{but in this case} \illegible{1} \unsure{universal}, should have served me as a guiding thread for my notations.}
\[
\begin{aligned}
(93)\qquad&\Sigma_{\rho,\sigma},\ B_{\rho,\sigma},\
\overset{\inserted{S_0G_{\rho,\sigma},}}{G_{\rho,\sigma}},\
\gamma_{\rho,\sigma},\ \deleted{\unsure{\Gamma_{\rho,\sigma}}},\
Z_\rho,\ Z_\sigma,\ \mathscr G_{\rho,\sigma},\ S\mathscr G_{\rho,\sigma}\\
&Z_\rho,\ Z_\sigma,\ S_0\Gamma_{\rho,\sigma}^{\deleted{\unsure{*}}0},\
\Gamma_{\rho,\sigma}^0,\ SZ_\rho,\ SZ_\sigma,\ S_0M_{\rho,\sigma}^0,\ M_{\rho,\sigma}^0.
\end{aligned}
\]

We put a prime $'$ everywhere to make them $\Sigma'_{\rho,\sigma}$, $B'_{\rho,\sigma}$ etc, and we remove the exponent $^0$ from $S_0\Gamma_{\rho,\sigma}^0$, $\Gamma_{\rho,\sigma}^0$, $S_0M_{\rho,\sigma}^0$, $M_{\rho,\sigma}^0$ (which is thus replaced by a $'$).

Other important special cases to \unsure{make explicit}:

$2^\circ)$ $\mathcal G=\mathrm{Sl}(2,\widehat{\mathbb Z})/\pm1$, $(\rho,\sigma)$ modular pair.
\[
\begin{aligned}
(94)\quad\left\{\begin{aligned}
B'_{\rho,\sigma}&\simeq
\left\{\begin{pmatrix}\mu&\lambda\\0&1\end{pmatrix}
\mid\mu\in\widehat{\mathbb Z}^{*},\ \mu\equiv1\ (3),\ \lambda\in\widehat{\mathbb Z}\right\}\\
&\inserted{L_0=\left\{\begin{pmatrix}1&\lambda\\0&1\end{pmatrix}\mid\lambda\in\widehat{\mathbb Z}\right\}}\\
&\deleted{\unsure{\mathrm{Gl}(\widehat{\mathbb Z})/\mu_2(\mathbb Z)}}\\
G'_{\rho,\sigma}&\simeq\mathrm{Gl}'(2,\widehat{\mathbb Z})/\pm1,\quad\text{where}\quad
\mathrm{Gl}'(2,\widehat{\mathbb Z})\overset{\mathrm{dfn}}=
\{u\in\mathrm{Gl}(2,\widehat{\mathbb Z})\mid\det u\equiv1\ (3)\}\\
&\deleted{\unsure{G\,\mathbb Z\,B/\mu_2}}\\
Z'_{\rho}&=\{c\rho+d\mid c^2+cd+d^2\in\widehat{\mathbb Z}^{*}\}/\pm1\\
&\inserted{\simeq\unsure{T_\rho|\mathcal T_\rho}(\widehat{\mathbb Z})/\mu_2(\mathbb Z)}\\
&\deleted{\text{this is the}\ \unsure{T_\rho}(\widehat{\mathbb Z})/\mu_2(\mathbb Z)\ \text{\unsure{with notations of} \S48}}\\
Z'_{\sigma}&=\deleted{\unsure{\{t\sigma+u\mid t^2+u^2\in\widehat{\mathbb Z}^{*}\}/\pm1}\quad\text{\unsure{subgroup} of \S48}}\\
&\inserted{N'_{\sigma}(\widehat{\mathbb Z})/\mu_2(\mathbb Z)}\quad\text{with notations of \S48}\\
S_0M'_{\rho,\sigma}&\simeq\deleted{\text{\unsure{is the}}}\
\unsure{SM''_{\rho,\sigma}|SM'_{\rho,\sigma}}(\widehat{\mathbb Z})/\Sigma\quad\text{with notations of \S48}\\
M'_{\rho,\sigma}&\simeq\deleted{\unsure{SM'_{\rho,\sigma}(\widehat{\mathbb Z})}}\
M'_{\rho,\sigma}(\widehat{\mathbb Z})/\Sigma\quad\text{id}\\
\Gamma'_{\rho,\sigma}&\simeq\mathcal T'_{\rho,\sigma}(\widehat{\mathbb Z})/\Sigma_0\quad\text{with notations of \S48}
\end{aligned}\right.
\end{aligned}
\]
\margin{this is $G(\widehat{\mathbb Z})/\mu_2(\mathbb Z)$ with notations of \S48}
\mspage{521}


\scanpage{523}

$3^\circ)$ Let us take $\mathcal G=\mathrm{Sl}(2,\widehat{\mathbb Z})$, with modular pair.
\[
(95)\quad\left\{\begin{aligned}
B'_{\rho,\sigma}&=
\left\{\begin{pmatrix}\mu&\lambda\\0&1\end{pmatrix}\mid
\mu\in\widehat{\mathbb Z}^{*},\ \mu\equiv1\ (3),\ \lambda\in\widehat{\mathbb Z}\right\}\\
G'_{\rho,\sigma}&=\mathrm{Gl}'(2,\widehat{\mathbb Z})
=\{u\in\mathrm{Gl}(2,\widehat{\mathbb Z})\mid\det u\equiv1\ (3))\\
Z'_\rho&=\{c\rho+d\mid
\overset{\inserted{c,d\in\widehat{\mathbb Z},\ \text{such that}}}{c^2+cd+d^2\in\widehat{\mathbb Z}^{*}}\}\\
Z'_\sigma&=T'_\sigma\cdot\{1,\tau_\infty\},\quad\text{where}\quad
\left\{\begin{aligned}
T'_\sigma&\overset{\mathrm{dfn}}=
\{t\sigma+u\mid t,u\in\widehat{\mathbb Z},\ t^2+u^2\in\widehat{\mathbb Z}^{*}\}\\
\tau_\infty&=\begin{pmatrix}-1&0\\0&1\end{pmatrix}
\end{aligned}\right.\\
S_0M'_{\rho,\sigma}&=S_0M'_{\rho,\sigma}(\widehat{\mathbb Z})/\Sigma_0\quad\text{with notations \S48}\\
M'_{\rho,\sigma}&=M'_{\rho,\sigma}(\widehat{\mathbb Z})/\Sigma_0\quad\rule{3cm}{0.3pt}\\
S_0\Gamma'_{\rho,\sigma}&=\unsure{S_0\mathcal T'_{\rho,\sigma}|S\mathcal T'_{\rho,\sigma}}(\widehat{\mathbb Z})/\Sigma_0\quad\rule{3cm}{0.3pt}\\
\Gamma_{\rho,\sigma}&=\mathcal T'_{\rho,\sigma}(\widehat{\mathbb Z})/\Sigma_0\quad\rule{3cm}{0.3pt}
\end{aligned}\right.
\]

\deleted{
$\overline{\text{\uline{III}}}$ \uline{The group} $\unsure{M_{\rho,\sigma}|M_\rho}$

Let us consider the homomorphisms

$(1)$
\diagram{\begin{tikzcd}[ampersand replacement=\&,column sep=large]
\deleted{\widehat{\mathbb Z}}
 \arrow[r,shift left=1.7em,"{\deleted{n\mapsto\rho^n}}"]
 \arrow[r,"{\deleted{n\mapsto\sigma^n}}"]
 \arrow[r,shift right=1.7em,"{\deleted{n\mapsto\epsilon_0^n}}"']
\&\deleted{\mathcal G}
\end{tikzcd}}
\deleted{\unsure{with image} $\unsure{L_0^{\mathcal G}}$}
\diagram{\begin{tikzcd}[ampersand replacement=\&,column sep=large]
\widehat{\mathbb Z}\arrow[r]\&\mathcal G\\
n\arrow[r,mapsto]\&\epsilon_0^n
\end{tikzcd}}
where $\epsilon_0=\sigma^{-1}\rho$
(with image $L_0^{\mathcal G}$)

and let $J$ be \deleted{the intersection of the} \inserted{\unsure{the}} kernels, in such a way that \deleted{\unsure{the image} (1), \unsure{factors}} \deleted{\unsure{through a}}
\[
(2)\qquad
\bigl(\widehat{\mathbb Z}^{\deleted{\unsure{*}}}/J\bigr)^*
\simeq\operatorname{Aut}(L_0^{\mathcal G}).
\]
We have \deleted{thus} an obvious homomorphism
\[
(3)\qquad
B'_{\rho,\sigma}\longrightarrow\operatorname{Aut}L_0^{\mathcal G}
\simeq(\widehat{\mathbb Z}/J)^*
\]
trivial on $L_0$, \deleted{\unsure{i.e. factors}} which gives a corresponding \unsure{homomorphism}
\[
(4)\qquad
\gamma'_{\rho,\sigma}\xrightarrow{\chi_\gamma}(\widehat{\mathbb Z}/J)^*,
\]
hence on composing with $G'_{\rho,\sigma}\longrightarrow\gamma'_{\rho,\sigma}$
\[
(4\text{bis})\qquad G'_\rho\xrightarrow{\chi_G}(\widehat{\mathbb Z}/J)^*.
\]
}
\mspage{522}


\scanpage{524}
\mspage{565}

$\overline{\text{\uline{III}}}$ \deleted{Let us consider the}

Let us consider the extension $\mathrm{Gl}(2,\mathbb Z)^\wedge$ of $\mathcal G_{0,3}^{\deleted{\unsure{+}}\wedge}$ by $\{\pm1\}$:

$(1)$
\diagram{\begin{tikzcd}[ampersand replacement=\&,column sep=1.4em,row sep=1.6em]
\&\&1\arrow[d]\&1\arrow[d]\&\\
1\arrow[r]\&\{\pm1\}\arrow[r,no head]\arrow[d,equal]\&
\mathrm{Sl}(2,\widehat{\mathbb Z})\arrow[r]\arrow[d]\&
\mathcal G_{0,3}^{+\wedge}\arrow[r]\arrow[d]\&1\\
1\arrow[r,no head]\&\{\pm1\}\arrow[r,no head]\&
\mathrm{Gl}(2,\widehat{\mathbb Z})\arrow[r,no head]\arrow[d]\&
\mathcal G_{0,3}^{\wedge}\arrow[r,no head]\arrow[d]\&1\\
\&\&\{\pm1\}\arrow[r]\arrow[d]\&\{\pm1\}\arrow[d]\&\\
\&\&1\&1\&
\end{tikzcd}}

We are interested in the group $\deleted{\unsure{M}}N_{1,1}^{\sim}$, defined \inserted{extension of $\deleted{\mathcal G_{0,3}^{\wedge}}M_{0,3}^{\sim}$ by $\{\pm1\}$,} as the inverse image of the extension $\mathrm{Gl}(2,\widehat{\mathbb Z})$ of $\mathrm{Gl}(2,\widehat{\mathbb Z})/\pm1$ by $\{\pm1\}$, via the canonical homomorphism
\[
(2)\qquad
M_{0,3}^{\sim}\deleted{\mathcal G_{0,3}^{\wedge}}
\longrightarrow\mathrm{Gl}(2,\widehat{\mathbb Z})'
=\mathrm{Gl}(2,\widehat{\mathbb Z})/\pm1
\]
(cf. \S45 \unsure{page}\ldots), so that we find a homomorphism of exact sequences,

\deleted{$\unsure{\tau_\infty(\epsilon_0)=\epsilon_0^{-1}}$}

$(3)$
\diagram{\begin{tikzcd}[ampersand replacement=\&,column sep=1.5em,row sep=2em]
1\arrow[r]\&\deleted{\unsure{M_{0,3}}}\{\pm1\}\arrow[r,no head]\arrow[d,equal]\&
\unsure{N_{1,1}^{\sim}|N_{0,3}^{\sim}}\arrow[r,no head]\arrow[d]\&M_{0,3}^{\sim}\arrow[r]\arrow[d]\&1\\
1\arrow[r]\&\{\pm1\}\arrow[r]\&\mathrm{Gl}(2,\widehat{\mathbb Z})\arrow[r]
\arrow[d,draw=none,"{\deleted{\unsure{\downarrow}}}"]\&
\mathrm{Gl}(2,\widehat{\mathbb Z})'\arrow[r,no head]
\arrow[d,draw=none,"{\deleted{\unsure{\downarrow}}}"]\&1\\
\&\&\deleted{\unsure{1}}\&\deleted{\unsure{1}}\&
\end{tikzcd}}

We find, \deleted{for $\unsure{\mathcal T_{1,1}^{\wedge}}$} \inserted{\unsure{a}} as the inverse image of $\deleted{\mathrm{Sl}(2,\mathbb Z)}\mathcal G_{0,3}^{+\wedge}\subset$ \deleted{in} $M_{0,3}^{\sim}$ in $N_{1,1}$, a $\mathrm{Gl}(2,\mathbb Z)^\wedge$, deduc-
\margin{\uline{N.B.} $\mathcal T_{1,1}\simeq\unsure{\mathrm{Gl}(2,\mathbb Z)_{\wedge}|\uline{\mathrm{Gl}(2,\mathbb Z)}}$.}
\mspage{523}


\scanpage{525}

ing a homomorphism
\[
(4)\qquad
\unsure{\mathrm{Sl}(2,\mathbb Z)^\wedge|\mathrm{Gl}(2,\mathbb Z)^\wedge}
\longrightarrow\deleted{\unsure{M}}N_{1,1}^{\sim}
\]
whose restriction to $\mathrm{Sl}(2,\mathbb Z)^\wedge$ fits \deleted{fits} into a canonical exact sequence

$(5)$
\diagram{\begin{tikzcd}[ampersand replacement=\&,column sep=1.8em]
1\arrow[r,no head]\&
\underbrace{\mathrm{Sl}(2,\mathbb Z)^{\unsure{\wedge}}}_{\mathcal T_{1,1}^{+\wedge}}\arrow[r]\&
N_{1,1}^{\sim}\arrow[r]\&
\underset{\substack{\Vert\\\Gamma_{0,3}^{\sim}}}{\Gamma_{\mathbb Q}^{\sim}}\arrow[r,no head]\&1
\end{tikzcd}}

which in turn fits into the diagram

$(6)$
\diagram{\begin{tikzcd}[ampersand replacement=\&,column sep=1.7em,row sep=2em]
1\arrow[r]\&\mathrm{Sl}(2,\mathbb Z)^\wedge\arrow[r]\arrow[d,"{\text{epimorphism}}"]\&
N_{1,1}^{\sim}\arrow[r,no head]\arrow[d,"{\text{epimorphism}}"]\&
\Gamma_{\mathbb Q}^{\sim}\arrow[r,no head]\arrow[d,"{\text{epimorphism}}","\chi"']\&1\\
1\arrow[r]\&\mathrm{Sl}(2,\widehat{\mathbb Z})\arrow[r,no head]\&
\mathrm{Gl}(2,\widehat{\mathbb Z})\arrow[r]\&\widehat{\mathbb Z}^{*}\arrow[r,no head]\&1
\end{tikzcd}}

We can write
\[
\begin{gathered}
(7)\qquad M_{0,3}^{\sim}\simeq
\overbrace{M_{0,3}^{\sim}(0)}^{\simeq\Gamma_{\mathbb Q}^{\sim}}
\cdot\overbrace{\mathcal G_{0,3}^{+\wedge}}^{\simeq\mathrm{Sl}(2,\mathbb Z)^{\wedge\prime}\simeq\mathrm{Sl}(2,\mathbb Z)^\wedge/\pm1}\\
\text{(semidirect)}
\end{gathered}
\]
in terms of \unsure{this} decomposition, and taking into account that we have a canonical homomorphism
\[
(8)\qquad M_{0,3}^{\sim}(0)\longrightarrow\mathrm{Gl}(2,\widehat{\mathbb Z})^{\unsure{\sharp|\prime}}
\]
\unsure{that|which} taking \unsure{the|there} \deleted{\unsure{homomorphism composed of}} \inserted{restriction of} (2) to $M_{0,3}^{\sim}(0)$, we find a canonical isomorphism
\mspage{524}

\scanpage{526}
\mspage{566}

\[
 (9)\qquad
 N_{1,1}^{\sim}\simeq
 \underbrace{N_{1,1}^{\sim}(0)}_{M_{0,3}^{\sim}(0)\simeq\Gamma_{\mathbb Q}^{\sim}}
 .\,\mathrm{Sl}(2,\mathbb Z)^\wedge .
\]

Thus we have an operation of
\inserted{$N_{1,1}^{\sim}(0)\simeq$}$\Gamma_{\mathbb Q}^{\sim}$
on $\mathrm{Sl}(2,\mathbb Z)^\wedge$
which I would like to make explicit, \deleted{We are going to}
\unsure{let us use|using} the \deleted{\unsure{isomorphism}} cartesian diagram

\[
(10)\qquad
\diagram{\begin{tikzcd}[ampersand replacement=\&,column sep=large,row sep=large]
\mathrm{Sl}(2,\unsure{\mathbb Z|\widehat{\mathbb Z}})^\wedge
 \arrow[r] \arrow[d]
 \arrow[dr,phantom,"\text{cartesian}"]
 \& \mathcal{G}_{0,3}^{+\wedge}
 \arrow[r,phantom,"\simeq"]
 \arrow[dr,"{\text{induced by (2)}}"']
 \& \mathrm{Sl}(2,\mathbb Z)^\wedge/\pm1
 \arrow[d,"{\unsure{q|q'}\ \text{canonical}}"]
 \\
\mathrm{Sl}(2,\widehat{\mathbb Z}) \arrow[rr]
 \& {} \& \mathrm{Sl}(2,\widehat{\mathbb Z})/\pm1
\end{tikzcd}}
\]

An element $\unsure{\widetilde{\nu}|\widetilde{\gamma}}$ of
$\Gamma_{\mathbb Q}^{\sim}$ corresponds to a pair
\deleted{$(\alpha,\beta)$ \unsure{of elements}}

\[
(11)\qquad
(\alpha,\beta)\in
\deleted{\unsure{\mathcal{G}_{0,3}}}
\underbrace{\widehat{\Pi}_{0,3}^{\deleted{\unsure{\wedge|\prime}}\prime}}_{
 \unsure{\{1,\tau_\infty\}}.\widehat{\Pi}_{0,3}}
\times\widehat{\Pi}_{0,3}
\]

satisfying \deleted{\unsure{of|to}} the loop equation

\[
(12)\qquad
\alpha\rho(\alpha)^{-1}=\beta\sigma(\beta)^{-1}\ell_1^\nu
\qquad (\nu\in\widehat{\mathbb Z})
\]

which operates on $\mathcal{G}_{0,3}^{+\wedge}\simeq
\mathrm{Sl}(2,\mathbb Z)^{\wedge\unsure{\prime}}$ by

\[
(13)\qquad
\left\{\begin{aligned}
u(\rho)&=\operatorname{int}(\alpha)\rho=\zeta\rho,
&\zeta&=\alpha\rho(\alpha)^{-1},\\
u(\sigma)&=\operatorname{int}(\beta)\sigma=\eta\sigma,
&\eta&=\beta\sigma(\beta)^{-1}
\end{aligned}\right.
\]

which is compatible with \deleted{an}\inserted{the} operation
\deleted{(\unsure{side} \illegible{2})}
\inserted{denoted $\unsure{\underline{u}|\underline{\alpha}}$}
on $\mathrm{Sl}(2,\widehat{\mathbb Z})$ induced
by $\operatorname{int}(\underline{u})$, where
$\underline{u}\in\mathrm{Gl}(2,\widehat{\mathbb Z})$ is the element

\mspage{525}

\scanpage{527}

defined \deleted{by} as the image of $u$ under (8), more
precisely

\[
(14)\qquad
\unsure{\widetilde{u}|\underline{u}}=
\unsure{i_{T_0}|i_{\Gamma_0}}(\mu),\qquad
\text{where }\mu=\chi(u)\text{ is the multiplier, and}
\]

\[
(15)\qquad
\left\{\begin{aligned}
\unsure{i_{T_0}|i_{\Gamma_0}} &:G_m\longrightarrow
\mathrm{Gl}(2,\widehat{\mathbb Z})\\
\mu&\longmapsto\begin{pmatrix}\mu&0\\0&1\end{pmatrix}.
\end{aligned}\right.
\]

We will have, more precisely

\[
(16)\qquad
\left\{\begin{aligned}
\underline{u}(\underline{\rho})&=\underline{\zeta}(\rho)
&\text{where }\underline{\zeta}&=\underline{\alpha}\rho(\underline{\alpha})^{-1}\\
\underline{u}(\underline{\sigma})&=\unsure{\underline{\eta}}(\sigma)
&\text{where }\underline{\eta}&=\varepsilon_{\unsure{\varphi|\rho}}
\underline{\beta}\sigma(\underline{\beta})^{-1}
\end{aligned}\right.
\]

\margin{the underlined ones denote the elements of
$\mathrm{Gl}(2,\widehat{\mathbb Z})$ ---
$\underline{\alpha}$ and $\underline{\beta}$
\unsure{are not determined} \unsure{except up to|except up to a} sign,
but $\underline{\zeta},\underline{\eta}$ are
\unsure{unambiguous|uniquely determined}.}

\deleted{$\varepsilon=\varepsilon_2(\mu)$}

with, as a reminder

\[
(17)\qquad
\varepsilon=\varepsilon_2(\mu)\in\pm1
\qquad
\left\{\begin{aligned}
\varepsilon&=+1&&\text{iff }\mu\equiv1\ (4)\\
\varepsilon&=-1&&\text{iff }\mu\equiv-1\ (4)
\end{aligned}\right.
\]

Of course, $\mathrm{Sl}(2,\mathbb Z)^\wedge$ is
\uline{generated} \inserted{\unsure{topologically}} by $\rho$ and $\sigma$,
and we identify $\pm1$ with the \unsure{central|center} subgroup
of $\mathrm{Sl}(2,\mathbb Z)$ hence of
$\mathrm{Sl}(2,\mathbb Z)^\wedge$, we will also write
$-g$ for the product of $g\in\mathrm{Sl}(2,\mathbb Z)^\wedge$
by $-1$. We will therefore have in $\mathrm{Sl}(2,\mathbb Z)^\wedge$

\[
(18)\qquad
\left\{\begin{aligned}
u(\rho)&=\deleted{\operatorname{int}\unsure{\zeta}(\rho)=}
\operatorname{int}(\widetilde{\alpha})(\rho)=\zeta\rho\\
u(\unsure{\sigma|\rho})&=
\varepsilon\operatorname{int}(\widetilde{\beta})(\sigma)=\eta\sigma
\end{aligned}\right.
\]
\unsure{and}

with
\[
(19)\qquad
\left\{\begin{aligned}
\zeta&=\widetilde{\alpha}\rho(\widetilde{\alpha})^{-1},\qquad
\eta=\varepsilon\widetilde{\beta}\sigma(\widetilde{\beta})^{-1}\\
&\varepsilon=\varepsilon_2(\mu)\in\{\pm1\}\subset
\mathrm{Sl}(2,\mathbb Z)^\wedge,
\end{aligned}\right.
\]

\mspage{526}

\scanpage{528}
\mspage{567}

where now $\widetilde{\alpha}$ and $\widetilde{\beta}$
denote arbitrary liftings of $\alpha$ and $\beta$ to
$\mathrm{Sl}(2,\mathbb Z)^\wedge$ --- they are determined
only up to sign, but
\inserted{$\operatorname{int}(\widetilde{\alpha}),
\operatorname{int}(\widetilde{\beta}),$}
$\zeta$ and $\eta$ are unambiguously determined by $u$.
\deleted{If \unsure{now}}
\inserted{Let us also note:}
\[
(20)\qquad u(\varepsilon_0)=\varepsilon_0^\mu
\quad\deleted{\text{\unsure{for}}}
\]

If now we have a \inserted{\unsure{surjective}} homomorphism
\[
(20)\qquad N_{1,1}^{\sim}\xrightarrow{\Psi}G
\]
inducing a homomorphism
\[
(21)\qquad
\mathrm{Sl}(2,\mathbb Z)^\wedge
\xrightarrow{\text{epimorphism}}
\underset{\substack{\Vert\\\operatorname{Im}\Psi}}{\mathcal{G}}
\subset G,
\]

the images of $\rho,\sigma$ in $\mathcal{G}$
(denoted by the same letters)
\unsure{satisfying|satisfy}
\[
(22)\qquad
\rho^6=\sigma^4=1,\qquad
\rho^3=\sigma^2\quad\text{and hence}
\]
\deleted{$[\rho^3,\sigma^2]=1$ \unsure{which suffices}}
\[
[\rho^3,\sigma]=[\sigma^2,\rho]=1
\quad\text{i.e. }\rho^3,\sigma^2
\]
\deleted{to \unsure{imply}}
are \uline{central}

\margin{thus in $\mathcal{G}$ there is a canonical
involutive central element, denoted $-1$ \unsure{. .}}

and we can, for $u\in N_{1,1}^{\sim}(0)$, corresponding to
\unsure{elements} $\widetilde{\alpha},\widetilde{\beta}$ of
$\mathrm{Sl}(2,\mathbb Z)^\wedge$
(determined modulo \uline{sign}), hence $\zeta,\eta$,
try to find an \inserted{\unsure{endomorphism|automorphism}}
$u_{\mathcal{G}}$ of $\mathcal{G}$,
satisfying the appropriate relations
\[
(23)\qquad
u_{\mathcal{G}}(\rho_{\mathcal{G}})
=\zeta_{\mathcal{G}}\rho_{\mathcal{G}},\qquad
u_{\mathcal{G}}(\sigma_{\mathcal{G}})
=\eta_{\mathcal{G}}\sigma_{\mathcal{G}},
\]
and bring into play \unsure{our|to the} $\widetilde{\alpha}$ and
$\widetilde{\beta}$. But we are

\mspage{527}

\scanpage{529}

obliged, since \deleted{\unsure{in general}} we do not
\inserted{\unsure{necessaril\textsuperscript{y}}}
$\widetilde{\alpha}\in\mathrm{Sl}(2,\mathbb Z)^\wedge$,
only $\widetilde{\alpha}\in\mathrm{Gl}(2,\mathbb Z)^\wedge$,
to bring in the element $\tau_\infty$ of
$\mathrm{Gl}(2,\mathbb Z)^\wedge$,
and \inserted{the action of} its image $\tau$ in $G$ on
$\mathcal{G}$. We will have, now taking

\[
(24)\qquad
\begin{aligned}
\overset{\deleted{\unsure{\mathcal{G}|\wedge}}}{\varepsilon_{0\mathcal{G}}}
&=\overset{\deleted{\unsure{\mathcal{G}|\wedge}}}{\sigma_{\mathcal{G}}}
\overset{\deleted{\unsure{\mathcal{G}|\wedge}}}{\rho_{\mathcal{G}}},
&\varepsilon_{1\mathcal{G}}&=\rho_{\mathcal{G}}\sigma_{\mathcal{G}}
\deleted{\unsure{=}}\,=\sigma_{\mathcal{G}}(\varepsilon_{0\mathcal{G}})
=\rho_{\mathcal{G}}(\varepsilon_{0\mathcal{G}})\\
&=-\sigma_{\mathcal{G}}^{-1}\rho_{\mathcal{G}}
&&=-\rho_{\mathcal{G}}\sigma_{\mathcal{G}}^{-1}
\end{aligned}
\]

\margin{which is not the usual convention
\unsure{up to} now\ldots}

the relations
\[
(25)\qquad
\tau(\sigma)=\sigma^{-1},\qquad
\tau(\rho)=\sigma(\rho^{-1})
\]

\margin{I drop the subscripts $\mathcal{G}$}

which imply, via (24) and (22)
\[
(26)\qquad
\tau(\varepsilon_0)=\varepsilon_0^{-1},\qquad
\tau(\varepsilon_1)=\varepsilon_1^{-1},
\]
and we introduce the group
\[
(27)\qquad
\deleted{\unsure{\{1,\tau\}}}
\mathcal{G}^{\unsure{\bullet|\prime}}
\overset{\mathrm{dfn}}{=}
\underset{\frac12\,\text{direct}}{\{1,\tau\}.\mathcal{G}},
\]
where $\{1,\tau\}$ operates on $\mathcal{G}$ by the formulas (25).

This said, we are in a position to paraphrase constructions
already made elsewhere.

\mspage{528}
